\section{Arithmetic sign tests and the remaining physical pairing}
\label{sec:outlook}
The exact remaining question is $S_N=G-D_N\ge0$. The positive
completion and the residual estimates do not settle it. This section
isolates an arithmetic cancellation concealed by the separate large
terms and gives a negative control with the same archimedean data.

\subsection{The prime discrepancy carries an essential signed term}
For $F=E_Lf$, put
\[
h_f(r)=\re\ip F{U_rF},\qquad
\psi_{\rm Ch}(x)=\sum_{p^m\le x}\log p,\qquad
E(r)=\psi_{\rm Ch}(e^r)-e^r+1.
\]
The function $h_f$ is continuous and zero for $r\ge L$; the locally
bounded-variation function $E$ satisfies $E(0)=0$.
Let $K_{\ge1}$ be the gamma tower with the lowest mass omitted.

\begin{proposition}\label{sign:discrepancy}
On the full logarithmic form domain,
\begin{equation}\label{sign:exact}
Q_L[f]=K_{\ge1}[F]+(4+w_0)\norm f^2
       -2\int_0^L e^{-r/2}h_f(r)\,dE(r).
\end{equation}
On the smooth compact-support core the final term can also be written
$2\int_0^L E(r)e^{-r/2}(h_f'(r)-h_f(r)/2)\,dr$.
\end{proposition}
\begin{proof}
The complete prime and pole terms are, respectively,
\[
-2\int_0^Le^{-r/2}h_f(r)\,d\psi_{\rm Ch}(e^r),\qquad
4\int_0^L\cosh(r/2)h_f(r)\,dr.
\]
Removing the prime density $e^r\,dr$ from the latter leaves
$2\int_0^Le^{-r/2}h_f(r)\,dr$, exactly the form of the kernel
$e^{-|x-y|/2}$. The lowest gamma summand is $4I$ minus that kernel.
This proves \eqref{sign:exact}, with every prime power included.
All manipulations apart from the gamma form are bounded at fixed
support, so the identity holds on the full domain. Stieltjes
integration by parts on the smooth core gives the derivative
version, with boundary terms zero at both endpoints.
\end{proof}
The derivative of a general logarithmic-domain autocorrelation is
not assumed to be a function. Formula~\eqref{sign:exact} is the
version used on that domain. The normalization agrees with the
Weil functional in \cite[Section 1.1]{Suzuki} and the digamma
special values in \cite{DLMF}.

\begin{theorem}[Smooth-density obstruction]\label{sign:density}
Replace the prime measure $d\psi_{\rm Ch}(e^r)$ by $e^r\,dr$,
retaining the exact gamma term, scalar contact and both poles.
The resulting closed semibounded form $Q_L^{\rm dens}$ has a
negative smooth compactly supported direction for every $L\ge2$.
At $L=2$, the unit $H_0^1(-1,1)$ input $f(x)=\cos(\pi x/2)$ obeys
\begin{equation}\label{sign:density-bound}
Q_2^{\rm dens}[f]<-2979/6125<0.
\end{equation}
\end{theorem}
\begin{proof}
Here $E=0$, so the exact operator is
$T_{\ge1,L}+(4+w_0)I$. For a zero extension in $H^1(\R)$,
\[
K_{\ge1}[F]\le
\left(\sum_{k\ge1}\frac2{a_k^3}\right)\norm{F'}^2
\le\frac{26}{125}\norm{F'}^2.
\]
The last inequality uses the first term at $a_1=5/2$ and its
 decreasing integral tail:
$2/a_1^3+1/(2a_1^2)=26/125$.
The digamma special value gives
$w_0=-\gamma_E-\pi/2-3\log2-\log\pi<-5$. Indeed
$\gamma_E>H_8-\log9>H_8-11/5=29/56>1/2$,
$\pi>3$, $\log2>2/3$, and $\log\pi>1$.
The logarithm bound follows from a positive exponential sum
 giving $e^{11/10}>3$; $e<3$ gives the last inequality.
These are elementary rational comparisons.
For the displayed unit cosine, $\norm{F'}^2=\pi^2/4$, and hence
\[
Q_2^{\rm dens}[f]<-1+\frac{26}{125}\frac{(22/7)^2}{4}
                  =-\frac{2979}{6125}.
\]
Approximation in $H_0^1$ gives a smooth compactly supported
negative input. Zero extension preserves its value at larger
support lengths.
\end{proof}
This control keeps the archimedean normalization and has positive
leading density $\psi_{\rm dens}(x)=x-1$. It also admits a generic
positive completion: choose $\alpha>-(4+w_0)$ and adjoin
$\sqrt{\alpha+4+w_0}\,I$ to the higher-mass gamma source.
Its Gram is $W_L^{\rm dens}+\alpha I$; compact resolvent and a
positive high sector are still available. A unit supersymmetric
cap can label these channels as well.
Thus the leading density and the general source architecture
cannot force the arithmetic ordering. This excludes dropping
 the discrepancy as a harmless error, not every method using
prime-counting estimates with additional signed information.

\subsection{Four mechanisms and what each would still have to prove}
The most concrete surviving approach is the mixed-response
certificate \eqref{mixed:certificate}, using the boundary columns
\eqref{boundary:column}. Its missing arithmetic lemma is a rule
for $K(L)$ and a proof that the low block covers the full
residual allowance at every required length. At prescribed finite
parameters this is a sufficient, independently checkable estimate;
the generic residual identity supplies no reason for its sign.
The noncompact boundary theorem rules out replacing its finite-input
tail estimate by an operator-norm truncation of the whole boundary
correction.

The ordered response gives a second approach. With $F_\ell$ as
above, a sufficient lemma is
\[
PW_LP-B^*F_\ell B\ge\theta^{\ell+1}\delta^{-1}B^*B.
\]
An approximation of $V$ needs its own propagated remainder from
\eqref{joint:words}. Individual prime translations must remain
between the compressed resolvents. Positivity of grouped response
terms is insufficient for their separately expanded words: for
$X=\operatorname{diag}(1,1/4)$ and
$Y=\left(\begin{smallmatrix}1&1\\1&1\end{smallmatrix}\right)\ge0$,
one has $\det(XY+YX)=-9/16$.
Only an estimate using the actual arithmetic entries can distinguish
these generic negative controls.

A third route would derive an arithmetic physical gluing law.
Retain the unused low output $Zp$ and define
\[
\mathcal E_Np=\sqrt\alpha\,(p,V^{1/2}H^{-1/2}Bp),
\qquad \mathcal E_N^*\mathcal E_N=D_N.
\]
An independently constructed contraction $\mathcal C_L$, obtained
for example as a projected unitary gluing, with
$\mathcal C_LZp=\mathcal E_Np$ would give $D_N\le G$; the
complementary output would retain the nonnegative difference.
Arithmetic must enter this intertwining identity, including its
contact normalization. Abstract existence of that contraction is
exactly the unresolved ordering: the map $Zp\mapsto\mathcal E_Np$
is contractive if and only if $p^*D_Np\le p^*Gp$.
The present unit cap changes none of these Grams and supplies no
such arithmetic identity.

The discrepancy formulation is a fourth route. The assertion
\[
-2\int_0^L e^{-r/2}h_f(r)\,dE(r)
\ge-(4+w_0)\norm f^2-K_{\ge1}[F]
\]
for every input is exactly Weil positivity rewritten, rather than
 a new sufficient principle. An independently approachable lemma
would need additional arithmetic control of these signed
 autocorrelations. Absolute bounds on $E$ alone have no favorable
sign against a derivative that changes sign. The smooth-density
control demonstrates that this discrepancy term carries essential
information.

\subsection{Arbitrary support does not require a compatible source limit}
At fixed $L,P$, the quantities
\[
M-\alpha I=PW_LP,\qquad B=QW_LP,\qquad H=QW_LQ|_{QX_L}
\]
are independent of the chosen positive reference shift.
Changing the permitted remote delay or adding inactive prime
channels therefore changes $G,D_N$ separately but not $S_N$.
Reference tuning cannot improve the remaining sign.

\begin{proposition}\label{sign:cofinal}
If $L_j\to\infty$ and, for every $j$, a valid finite cutoff and
an independent argument prove $S_{N_j,L_j}\ge0$, the complete
Weil form is nonnegative on every required compact support.
\end{proposition}
\begin{proof}
For any finite $L$ choose $L_j\ge L$. Zero extension to that
 envelope preserves the Fourier gamma form, both pole amplitudes,
and every active prime correlation. Additional return distances
have zero overlap. Apply \eqref{joint:criterion} at $L_j$.
\end{proof}
Thus a theorem proving the needed inequality separately for each
length suffices. The matrix dimensions, high-sector gaps and trial
sizes may depend on length. Neither a uniform positive gap for
 the full form nor a single all-prime physical system is required.
A fixed-support result alone cannot be extended upward; the
variational lowest Rayleigh infimum is nonincreasing under support
 enlargement.

Compatible physical sources across all envelopes are a separate
objective. The restriction law \eqref{joint:envelope} works inside
one fixed envelope; rebuilding changes the finite error. No such
all-envelope physical law is established here. The stationary
all-prime contact divergence retains its stated scope.
On the physical metric side, the chiral-equation identification
remains explicitly conditional. Its endpoint uniqueness theorem
and unit cap normalization do not determine the arithmetic sign.
The earlier scalar, ambient-holomorphic, finite-delay and finite
Schatten-class exclusions also retain their precise hypotheses.

\subsection{Reproducibility and review scope}
Ten small standard-library programs replay 59 labelled exact checks.
The ten newest check Green-kernel and boundary-column algebra,
ordered-tail counterexamples, mixed residual bounds, reference-shift
cancellation and rational constants for the smooth-density obstruction.
The previous 49 records are unchanged. None of these checks
certifies analytic domains, essential norm, Hodge theorems,
infinite-dimensional limits or the arithmetic matrix sign.

The package includes the numbered derivations, a manuscript coverage
map, build records, exact records and a replay validator. Complete
 dated snapshots preserve the earlier drafts, including the live
48-page version immediately before this sign analysis. Historical
claims have not been rewritten to contain later results.
Hashes establish artifact identity, not mathematical truth.

This is an AI-assisted working manuscript with internal review.
Independent specialist review remains necessary. Full Weil
positivity, an all-support positive physical pairing and the
Riemann hypothesis remain open.
